\scanpage{106}
\mspage{51}
% archive sheet number: 105
\unsure{Notons} \deleted{\unsure{les}}l'isomorphisme

\[
 \underline{\mathrm{Parf}}\,S^o\simeq(\underline{\mathrm{Parf}}\,S)^{\deleted{\unsure{o}}}
\]

(provenant du fait que $f:X\mathbin{\unsure{-|\to}}Y$ de $\mathrm{Cat}$ est une fibration parfaite ssi $f^o$ l'est). Cet isomorphisme \inserted{$X\longmapsto X^o$} est compatible avec les localiseurs (car $f\in W\Longleftrightarrow f^o\in W$), d'où

\[
 \mathrm{HOT}^{\mathrm{lc}}_W(S^o)\simeq\bigl(\mathrm{HOT}^{\mathrm{lc}}_W(S)\bigr)^{\deleted{\unsure{o}}}
\]

(isom. de catégories). On aura deux foncteurs canoniques entre les trois catégories de \inserted{$W$-}types d'homotopie relatifs sur $S$

\diagram{\begin{tikzcd}[ampersand replacement=\&,column sep=small,row sep=large]
 (\underline{\mathrm{Fibsc}}_0\,S)W_S^{-1}
 \arrow[r,"\approx"]\arrow[d,dashed,"i'"'] \&
 \mathrm{HOT}^{\mathrm{lc}}_W(S)\arrow[dl,"{i\ \deleted{\unsure{(\mathrm{adj.})}}}"']\arrow[dr,"{j\ \deleted{\unsure{\mathrm{adj.}}}}"] \&
 (\underline{\mathrm{Cofibsc}}_0\,S)W_S^{-1}\arrow[l,"\approx"']\arrow[d,dashed,"j'"]\\
 \mathrm{HOT}_W(S) \& \&
 \mathrm{HOT}^o_W(S)^{\deleted{\unsure{o}}}\overset{\mathrm{def}}{=}\mathrm{HOT}_W(S^o)
\end{tikzcd}}

Ces flèches \inserted{$i,j$} sont suggérées par les flèches en pointillés $i',j'$, provenant des inclusions

\[
 \underline{\mathrm{Fibsc}}_0\,S\subset\underline{\mathrm{Fibsc}}\,S,\quad
 \underline{\mathrm{Cofibsc}}_0\,S\subset\underline{\mathrm{Cofibsc}}\,S,
\]

qui définiraient $i,j$ à l'aide du choix de foncteurs quasi inverses pour les flèches horizontales. Mais $i,j$ ne seraient définis qu'à \unsure{isom.|isos} (\unsure{unique|uniques}) près. Il sera plus joli de \deleted{les définir} \inserted{choisir ces quasi-inverses} de façon canonique, à l'aide des flèches

\diagram{\begin{tikzcd}[ampersand replacement=\&,column sep=small]
 \underline{\mathrm{Fibsc}}_0\,S \& \underline{\mathrm{Parf}}\,S\arrow[l,"\Phi_0"']\arrow[r,"\Psi_0"] \& \underline{\mathrm{Cofibsc}}_0\,S
\end{tikzcd}}

$\left(\text{où }\begin{aligned}\Phi_0&=\Phi\mid\underline{\mathrm{Parf}}\,S\\ \Psi_0&=\Psi\mid\underline{\mathrm{Parf}}\,S\end{aligned}\right)$ du diagramme p.~8. La question principale

